% Lösungen Einheit 1 von 3 – Entscheidungsregel bestimmen: die Nullhypothese liefert p und die Binomialverteilung der Testgröße, die Alternative die Seite des Ablehnungsbereichs (rechtsseitig bei „mehr als“, linksseitig bei „weniger als“); die Grenze über kumulierte Wahrscheinlichkeiten so wählen, dass das Signifikanzniveau eingehalten wird (zusammenbau v0.1)
\einheitenkopf{Einheit 1 von 3 · Entscheidungsregel bestimmen: die Nullhypothese liefert p und die Binomialverteilung der Testgröße, die Alternative die Seite des Ablehnungsbereichs (rechtsseitig bei „mehr als“, linksseitig bei „weniger als“); die Grenze über kumulierte Wahrscheinlichkeiten so wählen, dass das Signifikanzniveau eingehalten wird}

\erg{8}{a) rechtsseitig \quad b) $k = 39$: $H_0$ wird abgelehnt bei mindestens 39 Kunden, sonst nicht; $P(X \ge 39) \approx 0{,}0393$ \quad c) $k = 17$: $H_0$ wird abgelehnt bei höchstens 17 Buchungen, sonst nicht; $P(X \le 17) \approx 0{,}0382$ \quad d) $k = 58$: $P(X \ge 58) \approx 0{,}0746 \le 0{,}08$, aber $P(X \ge 57) \approx 0{,}0975 > 0{,}08$ \quad e) Gezeigt ist nur, dass das Niveau bei 71 eingehalten wird; es fehlt, dass 71 die kleinste solche Grenze ist: $P(Y \ge 70) \approx 0{,}0575 > 0{,}05$, also gehört 70 nicht zum Ablehnungsbereich}
\erg{9}{a) bisher $P(Y \ge 18) \approx 0{,}0480$; bei 160 Teilen kleinstes $k$ mit $P(Y \ge k) \le 0{,}0480$: $k = 33$ ($P(Y \ge 33) \approx 0{,}0339$, $P(Y \ge 32) \approx 0{,}0524$) – nicht einfach die alte Grenze verdoppeln}
\erg{10}{a) Die Gegenhypothese lautet „weniger als 65\,\%“, der Test ist linksseitig: größtes $k$ mit $P(X \le k) \le 0{,}05$, also $k = 50$ ($P(X \le 50) \approx 0{,}0401$); abgelehnt wird bei höchstens 50 zufriedenen Kunden \quad b) Große Trefferzahlen sprechen für die Gegenhypothese; abgelehnt wird dort, wo das Ergebnis für sie spricht, also bei großen Trefferzahlen. \quad c) 29 liegt unter 31: $H_0$ wird nicht abgelehnt, die Maschine wird nicht gekauft; bewiesen ist damit nicht, dass höchstens 40\,\% kaufen – das Ergebnis reicht nur nicht zum Widerspruch}
